\section{Poincar\'e's shuffles: random walks on groups}

The second edition of \emph{Calcul des probabilit\'es}~\cite{poincare1912calcul} contains, in a chapter of miscellaneous questions added in 1912, a discussion of card shuffling that anticipates the modern theory of random walks on groups. Poincar\'e asked why repeated shuffles produce a deck in which every order is equally likely, and he answered by an argument from the theory of groups that refers to the work of Frobenius and Cartan; Mazliak~\cite{mazliak2012} describes it. The modern form is the following theorem.

\begin{theorem}[Convergence of random walks on finite groups]
\label{thm:random_walk_groups}
Let $G$ be a finite group and $Q$ a probability measure on $G$. The convolution powers $Q^{*n}$ converge to the uniform distribution $U$ on $G$ if and only if the support of $Q$ generates $G$ and is not contained in a coset of a proper normal subgroup. In that case, for every irreducible unitary representation $\rho\neq1$ the Fourier transform $\hat Q(\rho) = \sum_gQ(g)\rho(g)$ has spectral radius strictly less than one. Moreover, with $\|\mu-\nu\|_{TV} = \max_A|\mu(A)-\nu(A)|$,
\[
\|Q^{*n}-U\|_{TV}^2\leq\frac14\sum_{\rho\neq1}d_\rho\,\tr\big(\hat Q(\rho)^n(\hat Q(\rho)^n)^*\big),
\]
where the sum runs over the nontrivial irreducible unitary representations of $G$ up to equivalence and $d_\rho$ is the dimension of $\rho$.
\end{theorem}

\begin{proof}
If the support of $Q$ lies in a proper subgroup $H$, or in a coset $aN$ of a proper normal subgroup $N$, then $Q^{*n}$ gives mass one to $H$, or to the coset $a^nN$, which has $U$ mass at most $\frac12$. Then $\|Q^{*n}-U\|_{TV}\geq\frac12$ for every $n$, and there is no convergence. Now assume that the support generates $G$ and lies in no coset of a proper normal subgroup, and let $\rho\neq1$ be irreducible and unitary. Then $\hat Q(\rho)$ is a contraction. Suppose $\hat Q(\rho)v = \lambda v$ with $|\lambda| = 1$ and $|v| = 1$. Since $\lambda v$ is an average of the unit vectors $\rho(g)v$ for $g$ in the support, and the unit ball of a Hilbert space is strictly convex, $\rho(g)v = \lambda v$ for every $g$ in the support. Every element of the finite group $G$ is a product of elements of the support, so $v$ is an eigenvector of every $\rho(g)$. Irreducibility forces $\dim\rho = 1$, so $\rho$ is a nontrivial character $\chi$ with $\chi(g) = \lambda$ on the support, which then lies in a coset of the proper normal subgroup $\ker\chi$, contrary to assumption. Hence $\hat Q(\rho)$ has spectral radius less than one and $\hat Q(\rho)^n\to0$. The Fourier transform turns convolution into multiplication, $\widehat{Q^{*n}}(\rho) = \hat Q(\rho)^n$, and $\hat U(\rho) = 0$ for $\rho\neq1$. The displayed bound is the upper bound lemma of Diaconis and Shahshahani~\cite{diaconisshahshahani1981,diaconis1988}: Cauchy Schwarz turns total variation into an $L^2$ distance, and the Plancherel theorem for $G$ expresses the latter in terms of Fourier coefficients. Its right side is a finite sum of terms that tend to zero.
\end{proof}

The spectral radius cannot be replaced by the operator norm. On $S_3$ the measure uniform on $\{(12),(123)\}$ satisfies the conditions of the theorem, so it converges to uniform, but $\hat Q(\rho)$ has operator norm one on the two dimensional representation, with spectral radius $\frac12$. Poincar\'e did not compute rates, but the structure of his argument, a spectral gap on each nontrivial representation, is in our reading that of the theorem. Diaconis and Shahshahani showed in 1981 that random transpositions mix a deck of $n$ cards after about $\frac12n\log n$ steps~\cite{diaconisshahshahani1981}, and Bayer and Diaconis showed in 1992 that seven riffle shuffles suffice for a deck of fifty two~\cite{bayerdiaconis1992}. The method of arbitrary functions explains how a physical process forgets its initial condition; the random walk theorem explains how a sequence of independent randomizations forgets it. Both are theorems about convergence to a uniform law, and both are ancestors of the Markov chain Monte Carlo methods of Chapter~\ref{ch:spins}.

\section{The Maxwell, Poincar\'e and Borel lemma}

\begin{theorem}[Maxwell, Poincar\'e and Borel]
\label{thm:sphere_lemma}
Let $X^{(n)}$ be uniformly distributed on the sphere $\{x\in\R^n:|x|^2 = n\}$. For each fixed $k$ the law of the first $k$ coordinates $(X^{(n)}_1,\dots,X^{(n)}_k)$ converges in total variation to the standard Gaussian law on $\R^k$.
\end{theorem}

\begin{proof}
Let $Z$ be standard Gaussian in $\R^n$. By rotation invariance, $X^{(n)}\eqd\sqrt n\,Z/|Z|$. The first $k$ coordinates equal $Z_{1:k}\cdot\sqrt{n}/|Z|$, and $|Z|^2/n\to1$ almost surely by the law of large numbers. This gives convergence in law. For total variation, let $n>k$. The first $k$ coordinates have the density
\[
\begin{aligned}f_n(x) &= c_{n,k}\Big(1-\frac{|x|^2}{n}\Big)^{(n-k-2)/2},\quad |x|^2<n,\\
c_{n,k} &= \frac{\Gamma(n/2)}{(\pi n)^{k/2}\,\Gamma((n-k)/2)},\end{aligned}
\]
obtained by integrating over the sphere of radius $\sqrt{n-|x|^2}$ in $\R^{n-k}$, where $c_{n,k}$ is the constant that makes $f_n$ a probability density. By Stirling's formula $\Gamma(n/2)/\Gamma((n-k)/2)\sim(n/2)^{k/2}$, so $c_{n,k}\to(2\pi)^{-k/2}$, and $(1-|x|^2/n)^{(n-k-2)/2}\to e^{-|x|^2/2}$. Hence $f_n$ converges pointwise to the standard Gaussian density on $\R^k$. Since all these functions are probability densities, Scheff\'e's lemma upgrades pointwise convergence to convergence in total variation.
\end{proof}

The lemma is the microcanonical derivation of the Maxwellian: a gas with fixed total kinetic energy, uniformly distributed over the energy sphere, has Gaussian one particle velocities in the limit. It also underlies Kac's model in Chapter~\ref{ch:kinetic}, whose state space is exactly this sphere, and it is a relative of the concentration of measure phenomenon that L\'evy studied and that Milman later made a pillar of high dimensional geometry. The result is usually attributed to Poincar\'e, but Diaconis and Freedman could not find it in Poincar\'e's work~\cite{diaconisfreedman1987}. The earliest probabilistic statement they located is in Borel~\cite{borel1914}, who had studied the uniform measure on a ball of velocities in 1906~\cite{borel1906}, and they point to Maxwell for the idea. Their quantitative version gives a total variation bound of order $k/n$, the rate of the finite de Finetti theorem of Chapter~\ref{ch:axioms} for finitely many values.

\section{Poincar\'e on kinetic theory}

The paper of 1906, \emph{R\'eflexions sur la th\'eorie cin\'etique des gaz}~\cite{poincare1906}, examines the argument of Gibbs for the second law~\cite{gibbs1902} and distinguishes two hypotheses about probability. In the discontinuous hypothesis one has a very large number of similar systems, such as the molecules of a one dimensional gas, whose initial states are known exactly, and the probability of a region of phase space is the fraction of the systems that lie in it. In the continuous hypothesis the initial state of a single system is known only through a probability density. Under the first hypothesis the entropy integral is infinite, and the paper adopts the second. It also distinguishes the fine entropy, the limit of the cell entropies as the cells shrink, from the coarse entropy computed with cells so small that ordinary methods of observation cannot tell two states within a cell apart. It proves that the fine entropy is constant in time, so that any increase must come from coarse graining, and that, in the sign convention in which entropy increases, the coarse entropy is never smaller than the fine entropy. Gibbs had made the same point in the \emph{Elementary Principles}~\cite[Chapter~XII]{gibbs1902} with the picture of a coloring matter stirred through an incompressible liquid: its density measured on cells of fixed size tends to a uniform limit, although the density carried along with the liquid never changes.

These remarks are consistent with the structure of Lanford's theorem (Theorem~\ref{thm:lanford})~\cite{lanford1975}. The Boltzmann entropy that increases is the entropy of the one particle marginal, a coarse grained quantity; the Gibbs entropy of the full $N$ particle law is conserved; and the difference between them measures correlations created by collisions, which are invisible in the limit because they concentrate on a set of vanishing measure. In the language of Chapter~\ref{ch:axioms}, chaos is lost at finite $N$ and recovered in the limit, and Poincar\'e's two hypotheses correspond to the empirical measure of one gas and the de Finetti measure over gases.

\section{Poincar\'e and the theory of errors}

Poincar\'e's \emph{Calcul des probabilit\'es}~\cite{poincare1912calcul} also contains a careful discussion of Gauss's derivation of the normal law of errors. Gauss had shown in 1809~\cite{gauss1809} that if the arithmetic mean is always the most probable value of a location parameter, then the error law must be Gaussian. Poincar\'e examined the argument, noted that it presupposes a uniform prior on the location, and gave a version in which the prior is explicit and, with enough observations, does not affect the conclusion; Mazliak~\cite{mazliak2012} describes this analysis. It is an early instance of the observation that characterizations of the normal law by the optimality of the sample mean are statements about Bayes estimators under flat priors. In modern terms, the Gaussian is the only location family for which the sample mean is the maximum likelihood estimate of the location for every sample size, and under a flat prior the maximum likelihood estimate is the posterior mode. Teicher gave a rigorous proof that needs only continuity of the density and the property for samples of sizes two and three~\cite{teicher1961}. The three reasons for the normal law listed earlier in this chapter are thereby joined by a fourth, Gauss's: the Gaussian is the error law that makes averaging optimal.

\section{Poincar\'e and Hilbert}

Hilbert and Poincar\'e represented different philosophies of foundations. Hilbert's axiomatic method sought consistent formal systems from which theories could be deduced. Poincar\'e's conventionalism held that the axioms of geometry and mechanics are neither synthetic truths nor experimental facts but conventions, chosen for convenience and guided by experience; and he distrusted the formal turn in logic, which he satirized in his debates with Russell and Couturat. Yet on the sixth problem their views converge more than they diverge. Hilbert asked for the limiting process from atoms to continua. Poincar\'e had shown that the approach to equilibrium cannot hold for every initial condition, because of recurrence and reversibility, and that statements of this kind can hold for every smooth probability law on the initial conditions, as the method of arbitrary functions shows. The paper of 1906 makes the point for a very large number of small planets on circular orbits: the longitudes become uniformly distributed for every smooth law of the initial data, the recurrence time grows with the number of planets and is infinite in the limit, and an initial distribution that is not smooth, such as planets placed on a line, leaves a latent organization, a spiral wound too tightly to be seen, which a coarse observer takes for uniformity~\cite{poincare1906}. Boltzmann answered the objections of Loschmidt and Zermelo by moving from individual trajectories to probabilities. The modern answer to Hilbert, a limit theorem for the empirical measure under a probability law on initial conditions, as the number of particles grows and the time is held fixed, is due to Kac~\cite{kac1956} and, for hard spheres, Lanford~\cite{lanford1975} (Chapter~\ref{ch:kinetic}); it joins Hilbert's limit with Poincar\'e's probability measures.

The Bayesian reading adds one more convention to Poincar\'e's list: the prior on initial conditions. Poincar\'e had written that the probability of causes requires a hypothesis about prior probabilities; we add that the hypothesis can be checked, refined and computed with, and that the method of arbitrary functions explains when it washes out.

\section*{Exercises}
\begin{enumerate}
\item For the circle rotation $x\mapsto x+\alpha$ modulo one with $\alpha$ irrational, verify recurrence directly and compute the expected return time to an interval of length $\delta$ started uniformly in it.
\item Let $X$ be exponential with rate $\lambda$. Compute $\P(\{tX\}<\frac12)$ exactly and compare the error with the bound of Theorem~\ref{thm:arbitrary}.
\item Show that the Gaussian Poincar\'e inequality is an equality for the Hermite polynomial of degree one, and compute the ratio $\int|f'|^2\dd\gamma/\Var_\gamma(f)$ for the Hermite polynomial $\mathrm{He}_k$ of degree $k$ on $\R$.
\item Let $Q$ on $S_3$ give weight $\frac13$ to the identity and $\frac29$ to each of the three transpositions. Compute $\hat Q(\rho)$ for the two dimensional irreducible representation and for the sign representation, and find the rate of convergence to uniform. Explain why the walk that picks one of the three transpositions uniformly does not converge.
\item Derive the density of the first $k$ coordinates of a uniform point on the sphere of radius $\sqrt n$ in $\R^n$.
\end{enumerate}
